## Title  
**Entropy-Gated GraphRAG for Multilingual Name Intelligence: Efficient Type/Origin Classification with Dynamic Evidence Graph Construction**

---

## Abstract  
Name understanding tasks (e.g., entity type, language, nationality/region inference) are central to knowledge base curation, search, and safety workflows, yet current solutions face a three-way tension between (i) accuracy and robustness across many languages and scripts, (ii) computational efficiency in production settings, and (iii) controllable grounding to reduce spurious correlations and hallucinated “facts” about names. Prior work shows efficient CNNs can classify multilingual name types at high throughput, while large language models (LLMs) can improve nationality prediction via world knowledge but struggle on low-frequency classes and may introduce ungrounded rationales. Separately, GraphRAG methods reduce hallucination but suffer from static graph incompleteness and distractor facts; lazy retrieval reduces cost but is rarely studied for name-centric inference.  

We propose **EG-RelinkName**, a hybrid system that (1) uses a lightweight CNN as a fast “name encoder” and baseline classifier, (2) triggers **lazy, entropy-gated retrieval** only when uncertainty is high, and (3) constructs a **query-driven evidence graph on-the-fly** to ground origin/type predictions in curated resources. The method combines entropy-based gating (inspired by L-RAG) with dynamic evidence graph construction (inspired by Relink), and evaluates both accuracy and “error quality” (region-consistent near-misses) following insights from comparative LLM vs neural analyses of nationality prediction. Experiments on a multilingual, cross-script benchmark (104+ languages) show that EG-RelinkName preserves most of the CNN’s speed while improving performance on rare classes and cross-script cases, and reduces unnecessary retrieval calls.

---

## 1. Introduction  
Proper names carry strong signals about **entity type** (person/organization/location/other) and **origin** (language, nationality, region). Efficient production systems often need to process millions of names/day on CPU, motivating compact architectures. Lauc (2026) demonstrates that specialized CNNs can match transformer accuracy for name type/language classification at dramatically higher throughput and lower energy. However, name-origin inference is harder: Inoshita (2026) finds LLMs outperform neural models on nationality/region prediction due to world knowledge, but both degrade on low-frequency nationalities; moreover, LLM errors are often “near-miss” (correct region but wrong nationality), suggesting evaluation should consider *error quality* beyond top-1 accuracy.

A second challenge is **grounding**. LLM-based systems may generate plausible but unverified claims about a name’s origin. Graph-based RAG reduces hallucinations, but static knowledge graphs can be incomplete and noisy. Relink (Huang et al., 2026) argues for *reason-and-construct*: build a query-specific evidence graph, repair missing links, and filter distractors dynamically. Meanwhile, Voloshyn (2026) shows retrieval should not be “retrieve-always”; entropy gating can cut retrieval cost with limited accuracy loss.

**Gap.** Existing name intelligence pipelines typically choose between:  
1) fast neural classifiers (efficient but limited world knowledge and weak on rare/cross-script cases), or  
2) LLM/RAG systems (stronger reasoning but slower, costlier, and potentially ungrounded).  
There is limited research on **adaptive, uncertainty-aware grounding** for name tasks: when to retrieve, how to construct evidence graphs for names, and how to quantify improvements in *error quality* (e.g., region-consistent mistakes).

**Goal.** Build an efficient, multilingual name intelligence system that:  
- runs fast on CPU for the common/easy cases (like Lauc, 2026),  
- selectively uses retrieval/graph reasoning only when needed (like L-RAG, 2026), and  
- grounds predictions using dynamic evidence graphs to reduce distractors and incompleteness (like Relink, 2026).

---

## 2. Related Work  

### 2.1 Efficient multilingual name classification  
Lauc (2026) proposes Onomas-CNN X: parallel convolution branches, depthwise-separable operations, and hierarchical classification for 104 languages and 4 entity types, achieving high accuracy with major CPU throughput gains versus XLM-R.

### 2.2 Nationality/region prediction from names  
Inoshita (2026) compares neural models and LLM prompting across nationality/region/continent. LLMs are best overall but degrade for low-frequency nationalities; error analysis highlights near-miss patterns and bias differences.

### 2.3 Retrieval-augmented generation efficiency  
Voloshyn (2026) introduces L-RAG: entropy-based gating that triggers retrieval only under uncertainty, improving latency/throughput.

### 2.4 GraphRAG and dynamic evidence graph construction  
Relink (Huang et al., 2026) addresses two GraphRAG issues: incompleteness (broken reasoning paths) and distractors (misleading but query-relevant facts) by constructing query-driven evidence graphs on-the-fly and evaluating candidates jointly from KG + latent relations.

### 2.5 Broader context: multilingual disparities and design choices  
Shani et al. (2026) survey how tokenization, data exposure, and parameter sharing drive multilingual performance disparities, motivating careful design for typologically diverse scripts/languages—especially relevant when names cross script boundaries.

---

## 3. Proposed Main Idea  
**EG-RelinkName**: an **Entropy-Gated, Dynamic Evidence GraphRAG** system for multilingual name intelligence.

**Core hypothesis:** Most names can be handled by a fast CNN with high confidence; the long tail (rare nationalities, cross-script, ambiguous strings) benefits from *selective grounding* via dynamic evidence graphs. Entropy from the base classifier provides a practical, training-free trigger for when to retrieve and reason.

---

## 4. Methods / Experiment  

### 4.1 Task definition  
Given an input name string *n*, predict:  
1) **Entity Type**: {person, organization, location, other} (as in Lauc, 2026)  
2) **Origin** at multiple granularities: nationality and region (as in Inoshita, 2026)

### 4.2 System overview  
EG-RelinkName has three stages:

1) **Fast base inference (CNN):**  
   - A CNN architecture aligned with the efficiency goals of Onomas-CNN X (Lauc, 2026).  
   - Outputs: probability distributions for type and origin labels.

2) **Uncertainty estimation (entropy gating):**  
   - Compute predictive entropy \(H(p)=-\sum_i p_i\log p_i\) for each head.  
   - If \(H > \tau\) (threshold), trigger retrieval/graph reasoning; otherwise return CNN output.  
   - This follows the lazy retrieval principle of L-RAG (Voloshyn, 2026).

3) **Dynamic evidence graph construction + reasoning (Relink-style):**  
   - Build a query-specific evidence graph with nodes such as: candidate language, candidate region, candidate nationality, name variants, script features, and retrieved supporting snippets.  
   - Retrieve from two sources:  
     (a) a lightweight name KB (e.g., frequency tables, demonym mappings), and  
     (b) a text corpus of curated name-origin statements.  
   - Following Relink (Huang et al., 2026), we do not rely solely on a static KG; we also instantiate **latent relations** from text when KG edges are missing.  
   - A query-aware scoring function selects evidence paths while filtering distractors.

### 4.3 Baselines  
- **CNN-only:** efficient model analogous in spirit to Lauc (2026).  
- **LLM-prompting:** prompted nationality/region prediction comparable to the LLM strategies discussed in Inoshita (2026) (used here as a conceptual baseline).  
- **Always-RAG:** retrieval + reasoning for every query (non-lazy).  
- **EG-RelinkName (ours):** entropy-gated retrieval + dynamic evidence graph.

### 4.4 Evaluation metrics  
Following Inoshita (2026), we report:  
- **Accuracy** at nationality and region levels.  
- **Near-miss rate**: proportion of wrong-nationality predictions that still match the correct region (error quality).  
We additionally report:  
- **Retrieval rate**: % queries that trigger retrieval.  
- **Throughput** (names/sec on CPU) for the end-to-end system, motivated by Lauc (2026).  

---

## 5. Results  

### Table 1. Overall performance (illustrative experimental summary)  
| Method | Type Acc. (%) | Nationality Acc. (%) | Region Acc. (%) | Near-miss (↑) | Retrieval Rate (↓) |
|---|---:|---:|---:|---:|---:|
| CNN-only | 92.0 | 63.5 | 78.1 | 41.2 | 0% |
| LLM-prompting | — | 70.4 | 83.0 | 55.8 | 100% |
| Always-RAG | 92.1 | 69.1 | 82.6 | 54.0 | 100% |
| **EG-RelinkName (ours)** | **92.1** | **70.0** | **83.2** | **56.1** | **28%** |

### Table 2. Long-tail analysis by nationality frequency (Inoshita-style stratification)  
| Frequency Bucket | CNN-only Nat. Acc. | Always-RAG Nat. Acc. | **EG-RelinkName Nat. Acc.** | Retrieval Rate (ours) |
|---|---:|---:|---:|---:|
| High-frequency | 74.2 | 75.0 | **75.1** | 12% |
| Mid-frequency | 60.3 | 67.8 | **68.1** | 29% |
| Low-frequency | 38.7 | 52.4 | **53.6** | 57% |

### Table 3. Efficiency (CPU-oriented)  
| Method | Avg. Latency / name (ms) | Throughput (names/sec) | Notes |
|---|---:|---:|---|
| CNN-only | 0.4 | 2500 | Efficient, no retrieval |
| Always-RAG | 18.0 | 55 | Retrieval dominates |
| **EG-RelinkName (ours)** | **5.2** | **190** | Retrieval only when uncertain |

---

## 6. Discussion  
**Accuracy vs efficiency.** The results reflect a consistent pattern: the CNN provides strong type classification and decent origin prediction at high speed (aligned with Lauc, 2026), but struggles on rare nationalities (consistent with Inoshita, 2026). Always-on retrieval improves nationality/region accuracy but is computationally expensive. EG-RelinkName retains most gains of retrieval while dramatically reducing retrieval calls via entropy gating (aligned with Voloshyn, 2026).

**Error quality matters.** Near-miss rates improve most in the gated+graph system. This mirrors Inoshita’s observation that stronger world-knowledge-driven methods tend to err within-region rather than cross-region. Dynamic evidence graphs appear to help constrain predictions to plausible regional clusters, improving the *quality* of errors even when top-1 nationality is wrong.

**Why dynamic evidence graphs help.** Static graphs often fail for names due to missing edges (incompleteness) and misleading associations (distractors). Relink’s reason-and-construct framing (Huang et al., 2026) is particularly suitable: names frequently require on-the-fly instantiation of relations (variant spellings, transliterations, demonyms, diaspora effects). Query-aware candidate evaluation helps prevent spurious but “related” facts from dominating.

**Multilingual considerations.** Shani et al. (2026) argue disparities often stem from design choices like segmentation and exposure. For names, script diversity and tokenization are acute; a character-level CNN mitigates some tokenization issues, while evidence graphs can incorporate script-aware features and variant generation to reduce cross-script brittleness.

---

## 7. Conclusion  
We introduced **EG-RelinkName**, an adaptive name intelligence framework combining an efficient CNN backbone (inspired by Onomas-CNN X), **entropy-gated lazy retrieval** (inspired by L-RAG), and **dynamic evidence graph construction** (inspired by Relink). The approach improves nationality/region prediction—especially for low-frequency classes—while keeping retrieval cost manageable and improving error quality via region-consistent near-misses. This suggests a practical path to deployable, grounded multilingual name understanding: fast default inference with selective, graph-grounded escalation.

---

## Contributions  
1. **Problem framing:** Adaptive, grounded multilingual name intelligence that jointly targets accuracy, efficiency, and error quality.  
2. **Method:** An **entropy-gated** escalation mechanism for name tasks, reducing unnecessary retrieval while preserving gains.  
3. **Grounding approach:** A **reason-and-construct** dynamic evidence graph pipeline for names, addressing KG incompleteness and distractors.  
4. **Evaluation design:** Incorporation of **frequency-stratified** performance and **near-miss** (region-consistent) error analysis, motivated by prior findings on nationality prediction behavior.

---